\documentclass[Total.tex]{subfiles}

\begin{document}
	\textbf{Notions}
	
	\begin{itemize}
		\item For a set $E$, the cardinal and cardinal of collection of subsets are finite, denoted by
		
		\begin{gather*}
			|E|\qquad 2^{|E|}
		\end{gather*}
		
		\item Given a collection of sets $\mathcal{A}$, 
		
		\begin{gather*}
			min(\mathcal{A})\qquad max(\mathcal{A})
		\end{gather*}
		
		are well-defined.
	\end{itemize}
	
	\chapter{Definitions}
		\section{Independent Set and Circuits}
		
			\subsection{Matroid}
		
			\begin{Def}
				\label{Matroid-Matroid}(Whitney,1935) A matroid $M$ is an ordered pair $(E,\mathcal{I})$ consisting of a finite set $E$ and a collection $\mathcal{I}$ of subsets of $E$ having the following three properties:
				
				\begin{itemize}
					\item $\varnothing\in \mathcal{I}$;
					\item (Hereditary) If $I\in \mathcal{I}',I'\subseteq I$, then $I'\in \mathcal{I}$;
					\item (Independence Augmentation) If $I_1$ and $I_2$ are in $\mathcal{I}$ and $|I_1|<|I_2|$(which implies $I_2-I_1\not=\varnothing$), then there is an element $e$ of $I_2-I_1$ such that
					
					\begin{gather*}
						I_1\cup \{e\}\in \mathcal{I}
					\end{gather*}
				\end{itemize}
				
				The members of $\mathcal{I}$ are called \textbf{independent sets} of $M$. $E$ is called the \textbf{ground set}. 
				
				If $I\notin \mathcal{I}$, $I$ is dependent.
			\end{Def}
			
			
		
			
			\begin{proposition}
				Let $E$ be the set of column labels of $m\times n$ matrix over $F$. And let $\mathcal{I}$ be the set $X$ of subsets of $E$, whose vectors indexed by the $X$ is linear independent in $V(m,F)$, then 
				
				\begin{gather*}
					(E,\mathcal{I})
				\end{gather*}
				
				is a matroid.
			\end{proposition}
			
			\begin{proof}
				内容...
			\end{proof}
			
			\subsection{Circuits}
			
			\begin{Def}
				内容...
			\end{Def}
		\section{Bases}
			
			\begin{lemma}
				If $B_1,B_2$ are bases of Matroid $M$. Then
				
				\begin{gather*}
					|B_1|=|B_2|
				\end{gather*}
			\end{lemma}
			
			\begin{proof}
				内容...
			\end{proof}
			
			\begin{lemma}
				内容...
			\end{lemma}
			
			\begin{proof}
				内容...
			\end{proof}
			
			\begin{theorem}
				内容...
			\end{theorem}
			
			\begin{proof}
				内容...
			\end{proof}
			
			\begin{lemma}
				All the members of $\mathcal{B}$ have the same cardinality.
			\end{lemma}
			
			\begin{proof}
				内容...
			\end{proof}
		\section{Rank}
		\section{Closure}
		\section{Geometric Representation of Matroid of Small Rank}
		
	\chapter{Matroid Intersection}
	
		\section{Matroid Intersection Theorem}
			
			\begin{theorem}
				(Matroid Intersection Theorem, Edmond) Let 
				
				\begin{gather*}
					M_1=(S,I_1)\qquad M_2=(S,I_2)
				\end{gather*}
				
				be matroids, with rank functions $r_1,r_2$. Then the maximum size of set in $\mathcal{I}_1\cap \mathcal{I}_2$ is equal to 
				
				\begin{gather*}
					\min_{U\subseteq S} (r_1(U)+r_2(S-U))
				\end{gather*}
			\end{theorem}
			
			\begin{proof}
				内容...
			\end{proof}
	\chapter{Duality}
	\chapter{Minors}
	\chapter{Connectivity}
	\chapter{Representable Matroids}
	\chapter{Constructions}
	\chapter{Higer Connectivity}
	\chapter{Binary Matroid}
\end{document}